% New standalone increment. No prior source is changed or merged.
\documentclass[11pt]{article}
\usepackage[T1]{fontenc}
\usepackage[utf8]{inputenc}
\usepackage{lmodern,amsmath,amssymb,amsthm,mathtools,microtype}
\usepackage[margin=.88in]{geometry}
\usepackage{array,longtable,booktabs,xurl}
\usepackage[hidelinks]{hyperref}
\setlength{\emergencystretch}{3em}
\allowdisplaybreaks[2]
\numberwithin{equation}{section}
\newtheorem{theorem}{Theorem}[section]
\newtheorem{proposition}[theorem]{Proposition}
\newtheorem{lemma}[theorem]{Lemma}
\newtheorem{corollary}[theorem]{Corollary}
\theoremstyle{remark}\newtheorem{remark}[theorem]{Remark}
\newcommand{\E}{\mathbb E}
\newcommand{\dd}{\,d}
\newcommand{\one}{\mathbf1}
\newcommand{\KK}{\mathcal K}
\newcommand{\BB}{\mathcal B}
\newcommand{\Gp}{G_+}
\newcommand{\Clk}{\mathcal C_q}
\newcommand{\eps}{\varepsilon}
\newcommand{\Kfun}{\mathsf K}
\newcommand{\src}[2]{\textup{#1} \texttt{#2}}
\title{The response barrier and the completed weak passage}
\author{Standalone analytical increment for independent review}
\date{Version 1 --- 28 October 2026}
\begin{document}
\maketitle
\begin{abstract}
We exclude the lower response boundary in WPT's initialized
two-boundary theorem. A uniform core response of order $q^2$ is
retained with its sign. An all-label angle-difference estimate,
$|\theta-\psi|\le Cq(1+\log(1/q))$, makes the negative outside
response $o(q^2)$ under the already proved outside-mass guard.
The secondary response is at most $H/9+o(q^2)$. Consequently
$c\ge (\kappa_0/2)q^2\mu_{\KK}+H_{\Gp}/2>0$ on WPT's
closed stopped interval, with an explicit model constant
$\kappa_0=351/102400$. The first response boundary is therefore
$c=c_b$. A fixed-factor comparison of primary and clock amplitudes
also gives $H(t_b)\asymp c_b$ and
$t_b-t_0=[\log(c_b/H_0)+O(1)]/Q$. The clock defect itself remains
$o(1)$. No response-sign assumption is added to the initialized theorem.
\end{abstract}

\section{Domain, notation and the noncircular use of WPT}
\label{an03rbv1:domain}
Use the initialized balanced original flow and normalization of WPT:
\[
 P_1=\mathcal N(e_1,q^2I_2),\quad
 P_2=\mathcal N(\rho e_2,q^2I_2),\quad P=(P_1+P_2)/2,
 \quad \rho\in[13/20,7/10],\quad m(0,b)=q^{10}.
\]
Write $\lambda=\rho^2$, $Q=\lambda+q^2$, $\ell=\log(1/q)$,
$k=10(1-\lambda)$, and
\[
 U=\sqrt m\,a_\theta,\quad W=\sqrt m\,a_\psi,\quad
 R_\theta=\E_P[(X-f(X))X^T\one_{a_\theta\cdot X>0}].
\]
Angles are continuous lifts with equal initial values. For a fixed
cohort $S$, let $\dd\mu=m\dd\lambda_0$, $\mu_S=\int_S\dd\mu$,
\[
 c_S=Q^{-1}\E_{P_2}[X_2(f_S)_2],\quad
 E_S=\tfrac12\int_S(U_2^2+W_2^2)\dd\lambda_0,\quad
 J_S=\tfrac12\int_S(U_1^2+W_1^2)\dd\lambda_0.
\]
Thus $E_S+J_S=\mu_S$. Let
$z=(U_2+W_2)/2$, $a_-=(U_2-W_2)/2$,
$H_S=\int_Sz^2\dd\lambda_0$, and $A_S=\int_Sa_-^2\dd\lambda_0$.
The positive part notation for a scalar is $v_+=\max(v,0)$,
and $v_-=\max(-v,0)$; $a_-$ above is a signed coordinate,
not this positive-part operation.

Keep WPT's fixed partition, without reassigning any strip:
\begin{equation}
 \mathbb S^1=\KK\ \dot\cup\ \Gp\ \dot\cup\ \Sigma\ \dot\cup\ \BB,
 \qquad
 \BB=\mathcal C\ \dot\cup\ \mathcal Q_{3,\mathrm{int}}^0
       \ \dot\cup\ \mathcal S_\pi\ \dot\cup\ \mathcal S_{3\pi/2}.
 \label{an03rbv1:partition}
\end{equation}
Here $\mathcal C$ is the strong complement, $\mathcal D=\varnothing$,
and $\Clk\subset\Gp$ is the fixed clock cohort. Put
$H=H_{\Clk}$, $H_0=H(t_0)$,
$I(t)=Q\int_{t_0}^t(1-c)$, and $t_0=t_h$.

\src{WPT}{twoboundary} is an unconditional initialized theorem.
It provides $0<c(t_0)<c_b$ and all its estimates on
\begin{equation}
 [t_0,\tau],\qquad
 \tau=\min(t_-,t_+)<t_0+20\ell,\qquad
 t_-:=\inf\{t>t_0:c(t)=0\},\quad
 t_+:=\inf\{t>t_0:c(t)=c_b\}.
 \label{an03rbv1:twostop}
\end{equation}
In particular $0\le c\le c_b$ on this \emph{closed} interval.
Every application below of a source requiring this sign condition
is restricted to this interval, including a possible endpoint with
$c=0$. Positivity after $\tau$ is not used. The argument excludes
$\tau=t_-$ only after all its estimates have been obtained there.

We use WPT's constants and strict improvements, in particular
\begin{gather}
 \int_{t_0}^{\tau}r\le B,quad
 r=\sup_\theta|R_{11}(\theta)|,quad
 \|e_2\|_{P_1}+\|e_1\|_{P_2}\le M_{\rm res}q,quad
 |R_{12}|+|R_{21}|\le C(1+M_{\rm res})q,
 \label{an03rbv1:matrixinputs}\\
 E_{\Gp}<E_g=2c_b,quad J_G\le q^2K_Y^2,quad
 E_\Sigma\le C_\Sigma(q^5+q^{\Gamma_\nu}H),\quad
 \Gamma_\nu\ge\frac{289169}{924000}>\frac3{10},
 \label{an03rbv1:energyinputs}\\
 e\le3e_*/4\le3/16,\quad r_{\KK}=O(q^{1/20}),\quad
 |\theta|,|\psi|\le Rq\text{ on }\KK,\quad R=79/51,
 \quad \mu_{\BB}/q\le q^{1/10}.
 \label{an03rbv1:geometryinputs}
\end{gather}
Here $e$ bounds each coordinate error of the actual reference from
$(a,a)$, $r_{\KK}$ bounds each core coordinate's displacement from
that reference, and $a=\rho\Kfun(\rho a)$, with
$\Kfun(v)=\phi(v)+v\Phi(v)$. We have weakened some of WPT's
strict inequalities. Its improved geometry, not the raw inequality
$r_{\KK}<1$, will be essential for the core sign.

\section{Positive core response and the primary--secondary ledger}
\begin{lemma}[Explicit core sign and lower mass]
\label{an03rbv1:core}
Uniformly on \eqref{an03rbv1:twostop}, for sufficiently small $q$,
every core label has
\begin{equation}
 x=\theta/q\ge y_{\min},\quad y=\psi/q\ge y_{\min},\quad
 y_{\min}=\frac9{320},\qquad
 c_{\KK}\ge\kappa_0q^2\mu_{\KK},\quad
 \kappa_0=\frac{351}{102400}.
 \label{an03rbv1:corepositive}
\end{equation}
With the $d_q=\int_{t_0}^{\tau}|\eta|=o(1)$ of S2CP's core
mass equation,
\begin{equation}
 \mu_{\KK}(t)\ge\mu_{\KK}(t_0)e^{-d_q}
   =(1/2-o(1))e^{-d_q}\ge1/3.
 \label{an03rbv1:masslower}
\end{equation}
\end{lemma}
\begin{proof}
The elementary bound $\phi(0)>3/8$ suffices. For example,
$\pi/4=\int_0^1(1+x^2)^{-1}\dd x<1-1/3+1/5=13/15$
implies $\pi<32/9$. MOMC's positive root satisfies
$a\ge\rho\phi(0)\ge39/160$.
Choose $q$ last so that $r_{\KK}\le9/320$. Then
\[
 x,y\ge39/160-3/16-9/320=9/320.
\]
For the exact response, write $X=(qZ,V)$ on $P_2$, where
$V=\rho+qZ_2$ and $Z,Z_2$ are independent standard normals.
The trigonometric substitution used in \src{MOMC}{reduction} gives
\begin{align}
 c_{\KK}
 &=\frac{q^2}{Q}\int_{\KK} y\E[V(Z+xV)_+]\dd\mu
                  +O_R(q^4\mu_{\KK})\notag\\
 &=\frac{q^2}{Q}\int_{\KK}y\rho\Kfun(\rho x)\dd\mu
                  +O_R(q^4\mu_{\KK}).
 \label{an03rbv1:coremoment}
\end{align}
Indeed $\sin(qy)=qy+O_R(q^3)$ and
$a_\theta\cdot X=q\cos(qx)[Z+\tan(qx)V/q]$.
The compact-parameter Gaussian moment is smooth; Taylor expansion
in $V$ about $\rho$, with centered variance $q^2$, has error
$O_R(q^2)$. This is precisely the moment calculation in MOMC,
and is uniform over the current positive donor measure.
Since $x\ge0$, $\rho\Kfun(\rho x)\ge39/160$.
For $Q\le1$ the leading coefficient is at least $351/51200$;
the uniform $O_R(q^4)$ term is absorbed by half of it. This proves
\eqref{an03rbv1:corepositive} with a model-only constant.

S2CP \texttt{Bs}, as invoked on WPT's interval, gives
\[
 (\log\mu_{\KK})'=1-\mu_{\KK}+\eta,\quad
 \int|\eta|=d_q=o(1),\quad \mu_{\KK}(t_0)=1/2-o(1)<1.
\]
If $w(t)=\log\mu_{\KK}(t)<w(t_0)$, take the last crossing
of $w(t_0)$ before $t$ (or $t_0$ itself). Between that crossing
and $t$, $\mu_{\KK}<\mu_{\KK}(t_0)<1$, so $w'\ge\eta$.
Integration gives $w(t)\ge w(t_0)-d_q$; the other case is immediate.
This proves \eqref{an03rbv1:masslower}. In particular we have not
multiplied the pointwise overshoot bound by a logarithmic duration.
\end{proof}

\begin{lemma}[Signed primary and absorbed secondary charges]
\label{an03rbv1:explicit}
On the same interval,
\begin{align}
 c_{\Gp}&\ge\frac89H_{\Gp}
                  -C\mu_{\Gp}e^{-\rho^2/(4q^2)},
 \label{an03rbv1:primary}\\
 |c_\Sigma|&\le\frac19H+mathcal R_\Sigma(q),\qquad
 \mathcal R_\Sigma(q)
       \le C(q^5+q^{9/2}+q^{4+\Gamma_\nu})=o(q^2).
 \label{an03rbv1:secondary}
\end{align}
All constants in the second bound are fixed before $q$.
\end{lemma}
\begin{proof}
The exact additive decomposition is
\begin{equation}
 c=c_{\KK}+c_{\Gp}+c_\Sigma+c_{\BB}.
 \label{an03rbv1:ledger}
\end{equation}
LCRC \texttt{identity}, with its nonnegative negative-part
convention, says $c_S=H_S-A_S+\tau_S$.
LCRC \texttt{localclock/signgood}, already supplied by WPT,
gives $z>0$, $|a_-|\le z/3$, and $W_2=z-a_->0$ on $\Gp$.
Thus $A_{\Gp}\le H_{\Gp}/9$ and LCRC \texttt{taulower}
applies with no angular margin. This proves \eqref{an03rbv1:primary}.
Its tail is $o(q^2)$ because
$\mu_{\Gp}=E_{\Gp}+J_{\Gp}\le E_g+q^2K_Y^2$.

LCRC \texttt{responsecharge} holds for arbitrary signs. From
$J_\Sigma\le J_G\le q^2K_Y^2$ and
\eqref{an03rbv1:energyinputs} it yields
\[
 |c_\Sigma|\le C_\Sigma q^5+C_\Sigma q^{\Gamma_\nu}H
       +Cq^{9/2}+Cq^{2+\Gamma_\nu/2}\sqrt H.
\]
Young's inequality bounds the last term by
$H/18+Cq^{4+\Gamma_\nu}$; choose $q$ so that
$C_\Sigma q^{\Gamma_\nu}\le1/18$. The result is
\eqref{an03rbv1:secondary}. No bound of the form $o(H_0)$ is needed.
\end{proof}

\section{An all-label alignment bound with the required scope}
\label{an03rbv1:alignmentsection}
The fixed-$Cq$ alignment strip suggested in the route is stronger
than needed. We prove $Cq(1+\ell)$ globally instead. Its squared
cost times the outside mass is $o(q^2)$. This avoids imposing an
unproved invariant sign region on transitional strip labels.

\begin{lemma}[All-label alignment at entry]
\label{an03rbv1:entryalignment}
For every initial label, including the two Q3 endpoints,
\begin{equation}
 |\theta(t_0,b)-\psi(t_0,b)|\le C_{\rm al,0}q.
 \label{an03rbv1:deltaentry}
\end{equation}
The constant depends only on the proved S1 bounds and the model
box, not on LCRC's guards.
\end{lemma}
\begin{proof}
At the moving overlap $t_m=c_m\ell$,
$c_m=4/(1-\lambda)+1/2$, Q2EC \texttt{early} is an all-label
comparison and gives $|\theta-\psi|\le Cq^u$, where
$u=10-c_m\ge169/102>1$.
On $[t_m,t_0]$, MOMC \texttt{main} supplies the fixed right-half
family $\mathcal R$ in the chart $[-R,R]^2$, with mass
$M\le1/2$, remaining mass $\mu\le Cq^2\ell^2M$, and
$\int M\le3/2$. The last bound also follows from its mass law
by integrating $-\log(1-M)$ and its $o(1)$ integrated defect.

We need estimates for \emph{arbitrary receivers}, not an extension
of MOMC's receiver chart. Let
$T_s=\int_{\mathcal R}m a_\psi a_\theta^T\dd\lambda_0$ and
$A_i=\E_{P_i}[XX^T\one_{a_\theta\cdot X>0}]$.
Q2EC \texttt{fullgates}, valid at every receiver angle, says
\[
 R_\theta=\tfrac12[(I-T_s)A_1+A_2-B_2]+\mathcal R_\theta,
 \quad B_2=\E_{P_2}[f_{\mathcal R}X^T\one_{a_\theta\cdot X>0}],
 \quad |\mathcal R_\theta|\le C\mu+CM e^{-c/q^2}.
\]
Its donor moment bounds are
\[
 (T_s)_{11}=M+O(q^2M),\quad
 (T_s)_{12},(T_s)_{21}=O(qM),\quad (T_s)_{22}=O(q^2M),
\]
and $|(B_2)_{11}|+|(B_2)_{22}|\le Cq^2M$,
$|(B_2)_{12}|+|(B_2)_{21}|\le CqM$.
For all receiver gates, $(A_1)_{ii}\ge0$ and
$|(A_1)_{12}|\le\E_{P_1}|X_1X_2|\le Cq$.
Since $M\le1/2$, both diagonal coefficients of $I-T_s$ are
positive for small $q$. The two cross terms in its trace therefore
give at worst $-Cq^2M$. Also $\operatorname{tr}A_2\ge0$.
The target matrices are symmetric, so their antisymmetric part
vanishes exactly. These observations prove uniformly in the receiver
\begin{align}
 (\operatorname{tr}R_\theta)_-&\le
       C_R(q^2M+\mu+M e^{-c/q^2}),\notag\\
 |(R_\theta)_{12}-(R_\theta)_{21}|&\le
       C_R(qM+\mu+M e^{-c/q^2}).
 \label{an03rbv1:pretrace}
\end{align}
The donor gate-removal estimate is uniform because the receiver
indicator has absolute value at most one; the donor P1 margin
is at least $1/(2q)$. Thus no receiver-angle range was added.

For $\Delta=\theta-\psi$, direct subtraction of the exact angular
equations gives
\begin{equation}
 \Delta'=-\operatorname{tr}R_\theta\sin\Delta
                +(R_{12}-R_{21})\cos\Delta.
 \label{an03rbv1:deltaequation}
\end{equation}
While $|\Delta|<1$, the upper Dini derivative satisfies
\begin{equation}
 D^+|\Delta|\le(\operatorname{tr}R_\theta)_-|\Delta|
                         +|R_{12}-R_{21}|.
 \label{an03rbv1:deltacomparison}
\end{equation}
The negative-trace integral in \eqref{an03rbv1:pretrace} is
$O(q^2\ell^2)$, and its antisymmetric integral is
$O(q+q^2\ell^2)=O(q)$. Gronwall, starting with $O(q^u)$,
strictly improves the temporary $|\Delta|<1$ guard to $Cq$.
This proves \eqref{an03rbv1:deltaentry} without a P2-open-gate
assumption, a crossing-time assumption, or finite-noise cooperativity.
\end{proof}

\begin{lemma}[Alignment on the two-boundary interval]
\label{an03rbv1:postalignment}
There is a fixed $C_{\rm al}$ such that, for every label and
$t\in[t_0,\tau]$,
\begin{equation}
 |\theta(t,b)-\psi(t,b)|\le\delta_q:=C_{\rm al}q(1+\ell).
 \label{an03rbv1:deltapost}
\end{equation}
The constant may depend on $B,M_{\rm res},E_g,c_b$; its uses
below multiply only vanishing powers of $q$.
\end{lemma}
\begin{proof}
Put $N_-=\E_{P_2}[(f_2X_2)_-]$ and
$\mathcal G_b=5\mathrm e E_g/[9(1-c_b)]$.
The signed-gate proof of GESG yields the following useful direct
bound, not just a bound on the excess over $I'/2$:
\begin{equation}
 -R_{22}(\theta)
 \le\tfrac12M_{\rm res}q^2+	frac12\E_{P_2}[(f_2X_2)_+]
 =\tfrac12M_{\rm res}q^2+\tfrac12(Qc+N_-).
 \label{an03rbv1:negdiagonal}
\end{equation}
Indeed the P1 contribution costs at most
$\|e_2\|_{P_1}\|X_2\|_{P_1}/2$; in the negative of the P2
residual, discard the nonpositive selected $-X_2^2$ term and
then drop the gate from the positive part of $f_2X_2$.
The right side is nonnegative. S2CP \texttt{gate}, whose full
stopped hypotheses are supplied by WPT, gives
\[
 \int_{t_0}^{\tau}(Qc+N_-)
       \le\frac{10\mathrm e}{9}\frac{E_g}{1-c_b}+o(1).
\]
Therefore
\begin{equation}
 \int_{t_0}^{\tau}\sup_\theta(\operatorname{tr}R_\theta)_-
       \le B+\mathcal G_b+o(1).
 \label{an03rbv1:tracebudget}
\end{equation}
Here $(R_{11}+R_{22})_-\le|R_{11}|+(R_{22})_-$.
The $M_{\rm res}q^2$ integral is $o(1)$ on the duration $20\ell$.
The distinction between negative diagonal and absolute gate excess
is necessary in this step.

Let $a_*=C(1+M_{\rm res})$ be the uniform cross-entry constant
in \eqref{an03rbv1:matrixinputs}. Apply
\eqref{an03rbv1:deltacomparison} and Lemma
\ref{an03rbv1:entryalignment} while $|\Delta|<1$:
\[
 |\Delta(t)|\le e^{B+\mathcal G_b+1}
                 (C_{\rm al,0}q+20a_*q\ell).
\]
For example
$C_{\rm al}=e^{B+\mathcal G_b+1}(C_{\rm al,0}+20a_*)$
works. For sufficiently small $q$ this is less than $1/2$,
closing that guard. The estimate includes the closed endpoint
$\tau$ by continuity.
\end{proof}

\section{Signed outside response: a uniform moment estimate}
\begin{lemma}[Negative response of an almost aligned donor]
\label{an03rbv1:donorsign}
If $|\theta-\psi|\le\delta\le1/2$, its response per unit
radial mass has negative part at most
\begin{equation}
 C\{(q+\delta)^2+e^{-\rho^2/(4q^2)}\}.
 \label{an03rbv1:singlecharge}
\end{equation}
Consequently every fixed cohort $S$ on WPT's interval satisfies
\begin{equation}
 (c_S)_-\le C\{(q+\delta_q)^2+e^{-\rho^2/(4q^2)}\}\mu_S.
 \label{an03rbv1:cohortcharge}
\end{equation}
\end{lemma}
\begin{proof}
Put $s=\sin\theta$, $v=\sin\psi$. Stein's identity for the
two correlated Gaussian variables $X_2$ and $a_\theta\cdot X$
gives the exact formula
\begin{equation}
 A_q(s):=\E_{P_2}[X_2(a_\theta\cdot X)_+]
    =\rho q\Kfun(\rho s/q)+q^2s\Phi(\rho s/q).
 \label{an03rbv1:stein}
\end{equation}
The per-unit-mass response is $vA_q(s)/Q$, and $|v-s|\le\delta$.
Also, uniformly in $s$, the negative part of $A_q(s)$ can only
come from $X_2<0$, and hence is at most
$C e^{-\rho^2/(4q^2)}$ by the same fourth-moment/Chernoff
estimate as LCRC \texttt{taulower}.

If $v\ge0$, this proves the required estimate immediately.
If $v<0$ and $s\ge0$, then $s,|v|\le\delta$ and
$0\le A_q(s)\le C(q+s)$, since
$\Kfun(z)\le z+\phi(0)$ for $z\ge0$.
The negative response is at most $C\delta(q+\delta)$.
If $v<0$ and $s<0$, the second term in \eqref{an03rbv1:stein}
is nonpositive, so
\[
 (A_q(s))_+\le\rho q\Kfun(-\rho|s|/q),\qquad
 |v|\le|s|+\delta.
\]
The inequalities $\Kfun(-z)\le\phi(z)$ and
$\sup_{z\ge0}z\phi(z)<\infty$ give a cost
$C(q^2+q\delta)$. Division by $Q\ge169/400$ proves
\eqref{an03rbv1:singlecharge}. Integrate it against the positive
radial measure, using $(\int h)_-\le\int h_-$, to obtain
\eqref{an03rbv1:cohortcharge}.
This calculation remains valid at $s=0$ and $v=0$.
\end{proof}

\begin{proposition}[Outside sign budget]
\label{an03rbv1:outside}
Uniformly on $[t_0,\tau]$,
\begin{equation}
 (c_{\BB})_-\le C_{
m fixed}q^{31/10}(1+\ell)^2=o(q^2)
       \le\frac{\kappa_0}{4}q^2\mu_{\KK}.
 \label{an03rbv1:outsidebudget}
\end{equation}
In particular the proposed bound with the additional
$H_{\Gp}/9$ allowance holds.
\end{proposition}
\begin{proof}
Apply \eqref{an03rbv1:cohortcharge} with $S=\BB$ and the existing
guard $\mu_{\BB}\le q^{11/10}$. The exponential is smaller than
any positive power of $q$. Lemma \ref{an03rbv1:masslower} and
a final decrease of $q$ give the last inequality.

For completeness, the sharper component sources give the following
individual upper bounds, with the common harmless constant suppressed.
Set $L_q=1+\ell$, $s_q=O(\sqrt{\log\ell})$.
\begin{center}
\begin{tabular}{p{.25\textwidth}p{.27\textwidth}p{.34\textwidth}}
\toprule
Fixed ancestry & Supremum mass source & Negative response bound\\
\midrule
Strong complement $\mathcal C$ & S2CP \texttt{compmass},
$\mu_{\mathcal C}\le Cq^{1+\Delta_\nu}$,
$\Delta_\nu>3/25$ & $Cq^{3+\Delta_\nu}L_q^2$\\[3pt]
Original Q3 interior & GESG, $\mu\le Cq^{10}$ &
$Cq^{12}L_q^2$; also $Cq^{13}$ by its closed P2 gate\\[3pt]
Inner $\pi$ strip & PISB \texttt{main},
$\mu\le Cs_qq^{1+411/800}$ &
$Cs_qq^{3+411/800}L_q^2$\\[3pt]
$J_3$, including transit labels & LSB \texttt{main},
$\mu\le C(1+s_q)q^{1+769/4000}$ &
$C(1+s_q)q^{3+769/4000}L_q^2$\\[3pt]
Both strips' frozen zones & PISB/LSB,
$\mu\le Cq^{11}\sqrt\ell$ & $Cq^{13}\sqrt\ell L_q^2$\\
\bottomrule
\end{tabular}
\end{center}
For the interior's stronger optional estimate, GESG keeps the P2
gate margin at least $\rho\Lambda_q$, where
$e^{-k_G\Lambda_q^2}=q^3$, $k_G=169/3200$.
The fourth-moment/H\"older gate charge of a unit donor is $O(q^3)$;
multiplication by its $O(q^{10})$ mass gives $O(q^{13})$.
The generic alignment estimate already suffices there.
No comparison of the raw $\pi$-strip mass with $H_0$ is used.
Joined, transitional and frozen $J_3$ labels are all covered by
the same exact moment estimate; no whole-strip chart assertion is made.
\end{proof}

\section{The barrier, terminal amplitude and time}
\begin{theorem}[Strict response barrier on WPT's interval]
\label{an03rbv1:barrier}
For all sufficiently small $q$, with WPT's fixed constants,
\begin{equation}
 c(t)\ge\frac{\kappa_0}{2}q^2\mu_{\KK}(t)
                         +\frac12 H_{\Gp}(t)>0
 \qquad(t_0\le t\le\tau).
 \label{an03rbv1:barrierineq}
\end{equation}
Consequently $t_->t_+$ and $\tau=t_+=t_b<t_0+20\ell$.
In particular WPT's missing domain condition holds on
$[t_0,\min(t_+,t_0+20\ell)]$.
\end{theorem}
\begin{proof}
Combine the exact decomposition \eqref{an03rbv1:ledger},
Lemmas \ref{an03rbv1:core}--\ref{an03rbv1:explicit}, and
Proposition \ref{an03rbv1:outside}. Since $H\le H_{\Gp}$,
\[
 c\ge\kappa_0q^2\mu_{\KK}+\frac79H_{\Gp}
 -\mathcal R_\Sigma(q)-C\mu_{\Gp}e^{-\rho^2/(4q^2)}
 -(c_{\BB})_-.
\]
The outside charge uses at most $\kappa_0q^2\mu_{\KK}/4$.
The other two remainders are $o(q^2)$ uniformly, so by
$\mu_{\KK}\ge1/3$ they use at most another quarter for small $q$.
This proves the stronger bound with $7H_{\Gp}/9$ in place of
$H_{\Gp}/2$. If $\tau=t_-$, continuity gives $c(\tau)=0$,
contradicting \eqref{an03rbv1:barrierineq}. WPT already proved
finite $\tau$ before the duration boundary and excluded every
other guard exit. Thus $\tau=t_+$; positivity at this finite
endpoint and continuity also give $t_->t_+$.
\end{proof}

\begin{lemma}[Fixed primary-to-clock comparison]
\label{an03rbv1:clockcomparison}
There is a fixed $C_{
m ent}\ge1$ such that
\begin{equation}
 H_0\le H_{\Gp}(t_0)\le C_{
m ent}H_0,qquad
 H(t)\le H_{\Gp}(t)\le C_{
m ent}e^{2D_+}H(t),
 \quad D_+=C(C_0+q^2\ell),
 \label{an03rbv1:primaryclock}
\end{equation}
on the completed passage. In particular $H(t_b)\asymp c_b$.
\end{lemma}
\begin{proof}
We record the entry check; a whole-Q2 $q^5$ energy estimate would
not suffice here. The fixed primary group is
$c_0\ge K(C_0)q^\gamma$. Q2EC's proof of its entry theorem
keeps each of these labels in its controlled sector through $t_0$:
after its matched weak-axis band the actual angle is at least $qL$
and the target angle at least $16qL$, with fixed large $L$.
Q2EC \texttt{match, coordinatetransport} therefore compares its
positive $U_2(t_0)$ with $\widehat U_2(t_0)$ by fixed factors.
The matched band covers late crossers, including $b=\pi$; it
is not legitimate to take logarithms of their initial zero $U_2$.

For precision, Q2E \texttt{crossclock, crossamplitude} give
\[
 T_\times=\frac2\lambda\log\frac{c_0+q}{q(s_0+q)}+O(1),
 \qquad
 \widehat U_2(T_\times)\asymp q^5(c_0+q)/q.
\]
All such target crossings precede $t_0$ since
$4/\lambda\le1600/169<10$ and $\log A_q=O(1)$.
After crossing the exact coordinate rate is
\[
 (\log\widehat U_2)'=\lambda/2+d_q^{\rm tar}
       +\frac{\rho q}{2\sin\widehat\theta}
                         \Kfun(-\rho\sin\widehat\theta/q),
 \quad 0\le d_q^{\rm tar}\le q^2.
\]
Its extra integral is bounded uniformly while the endpoint angle
is at least $16qL$. To see the scope explicitly, the fixed weak-axis
layer has bounded duration; above $\pi/4$ the extra rate is
exponentially small. Below $\pi/4$, Q2EC's target version of the
positive cotangent rate gives
$-(\log\sin\widehat\theta)'\ge b/2$, $b=1/1000$.
Changing variable to $w=\sin\widehat\theta/q$ bounds the remaining
integral by $C\int_{cL}^{\infty}\Kfun(-\rho w)w^{-2}\dd w$.
Also $\int d_q^{\rm tar}=O(q^2\ell)$.
Thus, using $t_0=10\ell-\log A_q+o(1)$,
\begin{equation}
 \widehat U_2(t_0,b)^2
 \asymp q^{10}e^{\lambda t_0}(s_0+q)^2
 \asymp A_q^{-\lambda}q^k(s_0+q)^2
 \quad(b\in\Gp).
 \label{an03rbv1:primaryseed}
\end{equation}
All comparisons are fixed-factor ones. In particular at $b=\pi$
the regularized seed is $s_0+q=q$, not zero.
Q2EC's original entry signs give $z\le4U_2/3$.
Integrating the upper bound in \eqref{an03rbv1:primaryseed}
therefore yields $H_{\Gp}(t_0)\le C A_q^{-\lambda}q^k$.
Q2EC's clock seed gives the reverse comparison for $H_0$ on its
fixed interior subcohort. The upper and lower bounds on $A_q$
are supplied by MOMC. This proves the entry assertion.

RCET Theorem 2.1's relative-amplitude transport applies to both
positive-amplitude fixed subsets of $\Gp$, including the containing
set itself. Its required response range, duration, residual bounds,
$\int r\le B$, and entry signs/ratios all hold by WPT and Q2EC.
Its individual defect bound implies the aggregate bound on each
subset, giving \eqref{an03rbv1:primaryclock}; no disjointness is needed.

At $c(t_b)=c_b$, LCRC's absolute response charge, signgood and
$J_{\Gp}\le q^2K_Y^2$ give
\[
 c_{\Gp}(t_b)\le E_{\Gp}(t_b)+O(q^2)
                 \le\frac{10}{9}H_{\Gp}(t_b)+O(q^2).
\]
The core absolute response is $O(q^2)$, the secondary response
is $o(1)$ from its energy bound and $H\le O(c_b)$, and
$|c_{\BB}|\le C\mu_{\BB}=o(1)$. Hence
\begin{equation}
 \frac9{10}c_b-o(1)\le H_{\Gp}(t_b)
                  \le\frac98c_b+o(1).
 \label{an03rbv1:terminalprimary}
\end{equation}
Together with \eqref{an03rbv1:primaryclock}, this proves
$H(t_b)\asymp c_b$ with fixed constants, not a claimed limiting
ratio $H(t_b)/c_b$.
\end{proof}

\begin{theorem}[Initialized deep-regime weak passage]
\label{an03rbv1:completed}
For $\kappa=10$ and $\rho\in[13/20,7/10]$, choose the fixed
constants in WPT's order. For every sufficiently small $q$, the
original initialized flow has a first $t_b>t_0=t_h$ with $c(t_b)=c_b$,
and $t_b<t_0+20\ell$. On $[t_0,t_b]$ the barrier
\eqref{an03rbv1:barrierineq} and every conclusion of WPT
\texttt{twoboundary} hold. In particular:
\begin{align}
 \|e_1\|_{P_1}+\|e_2\|_{P_2}&\le C,&
 \|e_2\|_{P_1}+\|e_1\|_{P_2}&\le M_{\rm res}q/2,\notag\\
 J_G&\le q^2K_Y^2/4,&
 \int r&\le C_r\int|1-\kappa_1|+o(1)
                       \le C_r\log2+o(1),\notag\\
 H\le H_{\Gp}&\le9c_b/8+o(1),&
 E_{\Gp}&\le5c_b/4+o(1),\notag\\
 E_\Sigma&\le C_\Sigma(q^5+q^{\Gamma_\nu}H),&
 r_{\KK}&=O(q^{1/20}),\quad e\le3e_*/4,\notag\\
 (1/2-o(1))e^{-d_q}&\le\mu_{\KK}\le e^{d_q},& d_q&=o(1).
 \label{an03rbv1:conclusions}
\end{align}
The core stays in $|\theta|,|\psi|<Rq$, $R=79/51$.
The strong complement and the entire fixed outside population obey
\begin{align}
 \sup\mu_{\mathcal C}/q&\le Cq^{\Delta_\nu},\quad \Delta_\nu>3/25,
 \notag\\
 \sup\mu_{\BB}/q&\le C\{q^{\Delta_\nu}+q^9+s_qq^{411/800}
       +(1+s_q)q^{769/4000}+q^{10}\sqrt\ell\}=o(q^{1/10}).
 \label{an03rbv1:massconclusions}
\end{align}
Each corresponding time integral is at most $20\ell$ times the
displayed bound, hence tends to zero. All-endpoint S2 distance and
capped-gain conclusions remain valid. For CFA's canonical
core-only representation ($T=0$),
\begin{align}
 \int\eps_F&\le C_FC_{L_c}C_{\rm RCET}C_0E_g+o(1)<\bar F_F/2,
 \notag\\
 \int f_*&\le
 \frac{10\mathrm e E_g/9+C_V(E_g+K_Y\sqrt{E_g})}{Q(1-c_b)}
       +C_FC_{L_c}C_{\rm RCET}C_0E_g+o(1)
       \le\frac{e_*}{4C_{\rm pass}}.
 \label{an03rbv1:forcingconclusions}
\end{align}
The remainder norm includes value and induced $\ell^1$ receiver
derivative on $[-L_c,L_c]^2$. Its small fixed allowance is not
asserted to be $o(1)$ at fixed $c_b,C_0$.

The clock seed is $H_0\asymp A_q^{-\lambda}q^k$, with
$0<c_A\le A_q\le C_A$, and for every subinterval
\begin{equation}
 \log\frac{H(t)}{H(s)}=I(t)-I(s)+\delta(s,t),\qquad
 |\delta(s,t)|\le D_{\rm clk}=O(q^{1/8}+q^2\ell)=o(1).
 \label{an03rbv1:clock}
\end{equation}
At the endpoint $H(t_b)\asymp c_b$ and the exact time relation is
\begin{align}
 t_b-t_0
 &=\frac{\log(H(t_b)/H_0)-\delta(t_0,t_b)}{Q}
                                  +\int_{t_0}^{t_b}c,\notag\\
 0\le\int_{t_0}^{t_b}c
 &\le\frac{10\mathrm e}{9}\frac{E_g}{Q(1-c_b)}+o(1),\notag\\
 t_b-t_0
 &=\frac{\log(c_b/H_0)+O_{\rm fixed}(1)}Q
  =\frac{k\ell+\lambda\log A_q+\log c_b+O_{\rm fixed}(1)}Q.
 \label{an03rbv1:time}
\end{align}
The bounded phase error in the last line is not $o(1)$;
the exact subinterval clock defect in \eqref{an03rbv1:clock} is.
\end{theorem}
\begin{proof}
Theorem \ref{an03rbv1:barrier} identifies WPT's finite endpoint
as $t_b$, so its already proved conclusions hold without an
additional hypothesis. Lemma \ref{an03rbv1:clockcomparison}
supplies terminal comparability. Substitute it into WPT
\texttt{exacttime} and \texttt{cint}. The $1/Q$ in the integrated
response bound is retained. Since $c_b$ is fixed and small, that
integral is bounded independently of $q$. This proves
\eqref{an03rbv1:time}. No equality with denominator $Q(1-c_b)$
is inferred from a lower bound on $I'$.
\end{proof}

\section{Hypothesis audit, constants and S4 handoff}
\label{an03rbv1:audit}
The source statements used here have the following scopes. Their
stopped assumptions are conclusions of WPT on \eqref{an03rbv1:twostop},
not hypotheses imposed on a later interval.
\begin{longtable}{>{\raggedright\arraybackslash}p{.19\textwidth}
 >{\raggedright\arraybackslash}p{.45\textwidth}
 >{\raggedright\arraybackslash}p{.27\textwidth}}
\toprule
Source & Hypotheses needed here & Where supplied\\\midrule
WPT \texttt{twoboundary} & Original balanced isotropic initialization,
$\kappa=10$, ratio box, its fixed constant order, sufficiently small $q$.
No additional sign premise. & Model of Section 1; this source constructs
the interval with $0\le c\le c_b$.\\[4pt]
LCRC \texttt{ledger} & Fixed measurable cohorts, finite positive radial
mass; for \texttt{taulower}, $W_2\ge0$. No gate margin.
& Fixed partition; WPT signgood on $\Gp$.\\[4pt]
LCRC \texttt{localclock}; RCET/BULK & $0\le c\le c_b<1$, logarithmic
duration, bounded diagonal residuals, cross residual $O(q)$,
$\int r\le B$; positive entry $z$, $|a_-|/z\le1/4$,
entry ratio $qj/z\le C_0$ (clock $C_{\rm clk}q^{1/8}$),
small admissible $C_0$ and $q$. & WPT on its closed two-boundary
interval; Q2EC entry; unchanged constant order.\\[4pt]
S2CP \texttt{gate} & The preceding stopped residual/clock inputs,
$Y<K_Y$, $E_{\Gp}<E_g$, secondary joint guard, instantaneous core
chart/mass, outside supremum tending to zero, S1 entry and
$\mathcal D=\varnothing$. & WPT guards/improvements; no use of
the future passage forcing as a budget source.\\[4pt]
S2CP S2 and \texttt{Bs} & Same stopped guards; S2 additionally the
coarse c-form value/derivative integral and reference guard, SEP
and MOMC entry profiles; Bs uses core geometry, radial identity
and instantaneous outside bound. & WPT's matched joint bootstrap;
its canonical field, geometry, entry and outside improvements.\\[4pt]
MOMC \texttt{main}; Q2EC \texttt{early, fullgates} & Initialized
S1 flow on $[0,t_h]$, moving $t_m$, $M\le1/2$ on
$[t_m,t_h]$, strong donors on their fixed chart. Fullgates and
early are all-receiver/all-label statements. & Their unconditional
S1 conclusions; no weak-passage sign assumption.\\[4pt]
Q2E crossing; Q2EC matching/transport & Target-only Q2 closed sector
with regularized endpoints; original matched positive-coordinate
sector while angle $\ge qL$, fixed large $L$. & Q2EC's primary
threshold proof; $b=\pi$ handled by its band matching.\\[4pt]
GESG/PISB/LSB outside & Fixed ancestry cohorts, logarithmic stopped
duration, post-entry residual/clock guards and pre-entry cross
bound; no whole-strip lower-side or chart assumption. & WPT
\texttt{outside}; its use of Q3SB \texttt{crosspre}.\\[4pt]
CFA & Core-only positive normalized donor representation, stopped
geometry and residual/secondary budgets, paid outside/strip
integrals, clock-entry and forcing smallness. & Already matched
and proved in WPT; only its conclusions are inherited here.\\
\bottomrule
\end{longtable}

\paragraph{Constant order.}
Keep WPT's $\bar c=1/4$, $L_T=20$, $b_B=q^{1/10}$,
$R=79/51$, $R_{\rm LCRC}=2R$, $M_K=2$,
$B_s^0=\log2+1$, $L_c=R+1$, core cutoff $\chi=1/20$,
and admissible fixed $c_*,\nu$ first. Then fix the model passage
constants and the conservative, uniform $B,K_Y,M_{\rm res}$
(with $M_{\rm res}>2C_{\rm pre}$). Choose $\bar F_F$ and
small $c_b$, set $E_g=2c_b$, then the S2 distance/gain constants;
choose $C_0$ within LCRC/RCET and CFA's bounds; next fix
$K(C_0)$, SEP's entry constants and $C_\Sigma$; finally choose $q$.
This order is unchanged.

The new $y_{\min}=9/320$ and $\kappa_0=351/102400$ are model-only.
$C_{\rm al,0}$ is S1-only. $C_{\rm al}$ is explicitly allowed to
depend on $B,M_{\rm res},E_g,c_b$; it is used only in
$C_{\rm al}q(1+\ell)$ and in
$C_{\rm al}^2q^2(1+\ell)^2\mu_{\BB}$.
The secondary remainder constant depends on $C_\Sigma,K_Y$;
all its powers exceed two. $C_{\rm ent}$ depends on the fixed
primary threshold and matching constants, and terminal comparability
also uses the fixed RCET distortion. None of these new constants
chooses an earlier guard, radius or forcing allowance.

\paragraph{Endpoint and accuracy checks.}
The entry alignment proof includes every initial label, hence both
$b=\pi$ and $b=3\pi/2=-\pi/2$. The signed Stein bound allows
zero sine and negative sine without division by either.
The primary seed keeps $(s_0+q)$ and $(c_0+q)$ and uses matched
crossing before taking coordinate logarithms. There is no claim
that all late-joiner labels are charted or lower-side retained.
The proof uses the already improved core radius; the raw guard
$r_{\KK}<1$ alone would not prove positive core coordinates.
The energy alternative for outside response is not needed.

\begin{longtable}{>{\raggedright\arraybackslash}p{.25\textwidth}
 >{\raggedright\arraybackslash}p{.40\textwidth}
 >{\raggedright\arraybackslash}p{.26\textwidth}}
\toprule
S4 input & Certified at $t_b$ / through the passage & Accuracy\\\midrule
Response endpoint and phase & $c(t_b)=c_b$, strict positive barrier;
\eqref{an03rbv1:time}, $H(t_b)\asymp c_b$ & Exact endpoint;
bounded phase constant after replacing $H(t_b)$ by $c_b$\\[3pt]
Clock and seed & \eqref{an03rbv1:clock},
$H_0\asymp A_q^{-\lambda}q^k$ & $o(1)$ subinterval defect;
fixed-factor seed, not an exact coefficient\\[3pt]
Primary state/residuals & Signs, RCET ratio and relative-amplitude
transport, $J_G=O(q^2)$, \eqref{an03rbv1:conclusions} &
Fixed primary distortion; cross residual $O(q)$;
learning integral $\le\log2+o(1)$\\[3pt]
Retained core/reference & Fixed core, $r_{\KK}=O(q^{1/20})$,
$e\le3e_*/4$, $1/2-o(1)\le\mu_{\KK}\le1+o(1)$ &
Vanishing core radius; fixed reference allowance\\[3pt]
Secondary and outside & S2 entry-based energy, all-endpoint
distance/gain, \eqref{an03rbv1:massconclusions} and integrated
$\mu/q$ bounds & Explicit positive powers, with fixed constants\\[3pt]
Strong c-form/forcing & CFA representation and
\eqref{an03rbv1:forcingconclusions} & Small fixed integrated
value/derivative and forcing budgets\\
\bottomrule
\end{longtable}
These are inputs for a future weak return, matching or selection
argument. No continuation estimate beyond $t_b$ or S4 theorem
is proved here.

\begin{thebibliography}{20}\small
\bibitem{WPT} \path{AN03_weak_passage_theorem_v1.tex},
prefix \texttt{an03wptv1:}: \texttt{twoboundary, entry, guards,
improvements, exacttime, cint, gap}.
\bibitem{LCRC} \path{AN03_ledger_cap_residual_closure_v1.tex},
prefix \texttt{an03lcrcv1:}: \texttt{ledger, identity, taulower,
responsecharge, pointcap, localclock}.
\bibitem{S2CP} \path{AN03_S2_closure_package_v1.tex},
prefix \texttt{an03s2cpv1:}: \texttt{gate, S2, compmass, Bs}.
\bibitem{GESG} \path{AN03_global_envelope_signed_gate_v1.tex},
prefix \texttt{an03gesgv1:}: signed-gate proof, matrix values,
Q3 interior.
\bibitem{DECG} \path{AN03_distance_envelope_and_capped_gain_v1.tex},
prefix \texttt{an03decgv1:}: Proposition 3.2, Corollary 3.3.
\bibitem{PISB} \path{AN03_pi_strip_budget_v1.tex},
prefix \texttt{an03pisbv1:}: \texttt{main}, frozen zones.
\bibitem{LSB} \path{AN03_late_strip_budget_v1.tex},
prefix \texttt{an03lsbv1:}: \texttt{main}, delayed/frozen split.
\bibitem{CFA} \path{AN03_c_form_assembly_v1.tex},
prefix \texttt{an03cfav1:}: representation, assembly, forcing, passage.
\bibitem{MOMC} \path{AN02_moving_overlap_M_clock_passage_v1.tex},
prefix \texttt{an02momcv1:}: \texttt{main, reduction, massclaim,
outsideclaim}.
\bibitem{Q2EC} \path{AN02_Q2_entry_contracts_v1.tex},
prefix \texttt{an02q2ecv1:}: \texttt{early, matrixmoments, fullgates,
match, coordinatetransport}, primary-entry proof and clock seed.
\bibitem{Q2E} \path{AN02_Q2_boundary_layers_and_cohort_energy_v1.tex},
prefix \texttt{inc:AN02:q2completion:v1:}: \texttt{crossclock,
crossamplitude, U2rate}.
\bibitem{RCET} \path{AN03_residual_closure_and_eta0_transport_v1.tex},
prefix \texttt{an03rcetv1:}: Theorem 2.1, \texttt{transport}.
\bibitem{BULK} \path{AN03_bulk_clock_and_secondary_Q2_energy_v1.tex},
prefix \texttt{inc:AN03:bulkclock:v1:}: \texttt{crossinput,
pointseed}, Theorem 4.1.
\bibitem{PROFILE} \path{AN02_strong_entry_profile_and_tail_budget_v1.tex},
prefix \texttt{inc:AN02:profile:v1:}: Proposition 4.1 and Gaussian
gate-integral estimates.
\end{thebibliography}
\end{document}
